\documentclass[a4paper,12pt]{article}
\usepackage[utf8]{vietnam}
\usepackage[left=2cm,right=2cm,top=2cm,bottom=2cm]{geometry}
\usepackage{amsmath,amssymb}
\usepackage{ex_test}

\begin{document}

\begin{center}
%Tiêu đề
% Cách dùng: \tieude{tổng số trang}{mã đề}{tiêu đề riêng (để trống nếu dùng mặc định)}{role: hs hoặc gv}{lop: 10/11/12}
\newcommand{\tieude}[5]{
	\def\customtitle{#3}
	\ifthenelse{\equal{\customtitle}{}}
		{\def\thetitle{ÔN TẬP KIỂM TRA}}
		{\def\thetitle{\customtitle}}
	\noindent
	%Trái
	\begin{minipage}[b]{8.5cm}
		\centerline{\textbf{\fontsize{11}{0}\selectfont \hspace{2cm} TRƯỜNG THPT CHUYÊN HÙNG VƯƠNG}}
		\centerline{(\textit{Đề thi có #1\ trang})}
	\end{minipage}\hspace{1.5cm}
	%Phải
	\begin{minipage}[b]{9cm}
		\centerline{\textbf{\fontsize{11}{0}\selectfont \thetitle}}
		\centerline{\textbf{\fontsize{11}{0}\selectfont MÔN TOÁN, LỚP #5 (KHÔNG CHUYÊN)}}
	\end{minipage}\\
	\vspace{0.4cm}
	\ifthenelse{\equal{#4}{gv}}
		{\rightline{
			\setlength\fboxrule{1pt}
			\setlength\fboxsep{4pt}
			\fbox{\bf Mã đề #2}
		}}
		{}
	\vspace{4pt}
}
\end{center}

\begin{center}\fbox{\textbf{[THIẾU CÂU HỎI: chương 1 - dung\_sai\_cau\_lon]}}\end{center}

\begin{center}{\small\textit{Lỗi khi chạy hàm sinh câu (AttributeError): '\_cython\_3\_2\_8.cython\_function\_or\_method' object has no attribute 'randint'}}\end{center}

\begin{ex}%%[?]
Trong các câu sau, câu nào là mệnh đề?
\choice{Bác Hồ đã đi qua bao nhiêu quốc gia trong suốt hành trình 30 năm bôn ba tìm đường cứu nước?}
{Ai yêu nhi đồng bằng bác Hồ Chí Minh?}
{\True Mặt Trời là một hành tinh trong Hệ Mặt Trời}
{Bạn bao nhiêu tuổi?}


\loigiai{
Mệnh đề là câu khẳng định có tính đúng hoặc sai. Các câu hỏi, câu cảm thán, câu cầu khiến không phải là mệnh đề.
}
\end{ex}

\begin{center}\fbox{\textbf{[THIẾU CÂU HỎI: L10\_C1\_B2\_NB017 - trac\_nghiem]}}\end{center}

\begin{center}{\small\textit{Lỗi khi chạy hàm sinh câu (AttributeError): '\_cython\_3\_2\_8.cython\_function\_or\_method' object has no attribute 'randint'}}\end{center}

\begin{ex}%%[?]
Phát biểu bằng lời mệnh đề:
``$\forall n\in \mathbb{N}, n^2 < n$''.
\choice{Tồn tại số tự nhiên mà bình phương của nó lớn hơn hoặc bằng chính nó}
{Mọi số tự nhiên đều có bình phương bằng chính nó}
{\True Mọi số tự nhiên đều có bình phương nhỏ hơn chính nó}
{Có ít nhất một số tự nhiên có bình phương nhỏ hơn chính nó}


\loigiai{
Kí hiệu $\forall$ phát biểu là "mọi",
kí hiệu $\exists$ phát biểu là "tồn tại".
}
\end{ex}

\begin{ex}%%[?]
Cho hai mệnh đề sau:\\
            $P \colon$ ``$a$ là số chính phương'';\\
            $Q \colon$ ``Căn bậc hai của số $a$ là số nguyên''.\\
            Hãy phát biểu mệnh đề $P \Leftrightarrow Q$.
\choice{Nếu $a$ là số chính phương thì căn bậc hai của số $a$ là số nguyên}
{Vì $a$ là số chính phương nên căn bậc hai của số $a$ là số nguyên}
{\True Căn bậc hai của số $a$ là số nguyên nếu và chỉ nếu $a$ là số chính phương}
{Nếu căn bậc hai của số $a$ là số nguyên thì $a$ là số chính phương}


\loigiai{
Mệnh đề tương đương $P \Leftrightarrow Q$ được phát biểu dưới dạng ``$P$ khi và chỉ khi $Q$''.
}
\end{ex}

\begin{ex}%%[?]
Trong các khẳng định sau, khẳng định nào \textbf{sai}?
\choice{$\exists n\in\mathbb Z,\ n>282$}
{\True $\forall n\in\mathbb Z,\ n>282$}
{$382>282$}
{$283>282$}


\loigiai{
Khẳng định đúng là phương án đã chọn.
}
\end{ex}

\begin{center}\fbox{\textbf{[THIẾU CÂU HỎI: L10\_C1\_B2\_VD020 - trac\_nghiem]}}\end{center}

\begin{center}{\small\textit{Lỗi khi chạy hàm sinh câu (AttributeError): '\_cython\_3\_2\_8.cython\_function\_or\_method' object has no attribute 'choice'}}\end{center}

\begin{center}\fbox{\textbf{[THIẾU CÂU HỎI: L10\_C1\_B2\_VD021 - tra\_loi\_ngan]}}\end{center}

\begin{center}{\small\textit{L10\_C1\_B2\_VD021 (SA): không có Generator nào khớp trong Mapping.}}\end{center}

\begin{center}\fbox{\textbf{[THIẾU CÂU HỎI: L10\_C1\_B2\_VD020 - tra\_loi\_ngan]}}\end{center}

\begin{center}{\small\textit{L10\_C1\_B2\_VD020 (SA): không có Generator nào khớp trong Mapping.}}\end{center}

\begin{ex}%%[?]
Tại một sự kiện có 100 người tham gia khảo sát về hai sản phẩm A và B. Biết có 52 người chọn sản phẩm A, 50 người chọn sản phẩm B và 32 người chọn cả hai sản phẩm. Tính: 
\begin{listEX}[1]
\item Số người đã chọn ít nhất một sản phẩm? \SA[4]{$70$}
\item Số người không chọn sản phẩm nào? \SA[4]{$30$}
\end{listEX}

 \loigiai{
\begin{listEX}[1]
\item Theo nguyên lý bù trừ, số người chọn ít nhất một sản phẩm là: 52 + 50 - 32 = 70 (người).
\item Số người không chọn sản phẩm nào là: 100 - 70 = 30 (người).
\end{listEX}

}
\end{ex}

\begin{center}\fbox{\textbf{[THIẾU CÂU HỎI: L10\_C1\_B2\_VD021 - tu\_luan]}}\end{center}

\begin{center}{\small\textit{L10\_C1\_B2\_VD021 (TL): không có Generator nào khớp trong Mapping.}}\end{center}

\begin{center}\fbox{\textbf{[THIẾU CÂU HỎI: L10\_C1\_B2\_VD021 - tu\_luan]}}\end{center}

\begin{center}{\small\textit{L10\_C1\_B2\_VD021 (TL): không có Generator nào khớp trong Mapping.}}\end{center}



\end{document}